\documentclass[10pt,a4paper]{amsart}
\usepackage{amsmath,amssymb,amsthm,graphicx}
\usepackage{verbatim}
\newenvironment{reference}{\comment}{\endcomment}
%\newenvironment{reference}{}{}
\newenvironment{slogan}{\comment}{\endcomment}
\newenvironment{history}{\comment}{\endcomment}

% For commutative diagrams we use Xy-pic
\usepackage[all]{xy}

% We use 2cell for 2-commutative diagrams.
\xyoption{2cell}
% Theorem environments.
%
\theoremstyle{plain}
\newtheorem{theorem}[subsection]{Theorem}
\newtheorem{proposition}[subsection]{Proposition}
\newtheorem{lemma}[subsection]{Lemma}

\theoremstyle{definition}
\newtheorem{definition}[subsection]{Definition}
\newtheorem{example}[subsection]{Example}
\newtheorem{exercise}[subsection]{Exercise}
\newtheorem{situation}[subsection]{Situation}

\theoremstyle{remark}
\newtheorem{remark}[subsection]{Remark}
\newtheorem{remarks}[subsection]{Remarks}

\numberwithin{equation}{subsection}

% Macros
%
\def\lim{\mathop{\mathrm{lim}}\nolimits}
\def\colim{\mathop{\mathrm{colim}}\nolimits}
\def\Spec{\mathop{\mathrm{Spec}}}
\def\Hom{\mathop{\mathrm{Hom}}\nolimits}
\def\Ext{\mathop{\mathrm{Ext}}\nolimits}
\def\SheafHom{\mathop{\mathcal{H}\!\mathit{om}}\nolimits}
\def\SheafExt{\mathop{\mathcal{E}\!\mathit{xt}}\nolimits}
\def\Sch{\mathit{Sch}}
\def\Mor{\mathop{\mathrm{Mor}}\nolimits}
\def\Ob{\mathop{\mathrm{Ob}}\nolimits}
\def\Sh{\mathop{\mathit{Sh}}\nolimits}
\def\NL{\mathop{N\!L}\nolimits}
\def\CH{\mathop{\mathrm{CH}}\nolimits}
\def\proetale{{pro\text{-}\acute{e}tale}}
\def\etale{{\acute{e}tale}}
\def\QCoh{\mathit{QCoh}}
\def\Ker{\mathop{\mathrm{Ker}}}
\def\Im{\mathop{\mathrm{Im}}}
\def\Coker{\mathop{\mathrm{Coker}}}
\def\Coim{\mathop{\mathrm{Coim}}}

% Boxtimes
%
\DeclareMathSymbol{\boxtimes}{\mathbin}{AMSa}{"02}

%
% Macros for moduli stacks/spaces
%
\def\QCohstack{\mathcal{QC}\!\mathit{oh}}
\def\Cohstack{\mathcal{C}\!\mathit{oh}}
\def\Spacesstack{\mathcal{S}\!\mathit{paces}}
\def\Quotfunctor{\mathrm{Quot}}
\def\Hilbfunctor{\mathrm{Hilb}}
\def\Curvesstack{\mathcal{C}\!\mathit{urves}}
\def\Polarizedstack{\mathcal{P}\!\mathit{olarized}}
\def\Complexesstack{\mathcal{C}\!\mathit{omplexes}}
% \Pic is the operator that assigns to X its picard group, usage \Pic(X)
% \Picardstack_{X/B} denotes the Picard stack of X over B
% \Picardfunctor_{X/B} denotes the Picard functor of X over B
\def\Pic{\mathop{\mathrm{Pic}}\nolimits}
\def\Picardstack{\mathcal{P}\!\mathit{ic}}
\def\Picardfunctor{\mathrm{Pic}}
\def\Deformationcategory{\mathcal{D}\!\mathit{ef}}
\usepackage{cleveref}
\crefname{originalequation}{original equation}{original equations}
\Crefname{originalequation}{Original equation}{Original equations}
\makeatletter
\newlabel{equation-functorial-omega}{{00RQ}{source}{}{sourcepointer-1}{}}
\newlabel{equation-functorial-omega@cref}{{[originalequation][0][]00RQ}{[1][0][]source}{}{}{}}
\makeatother
\crefname{originalequation}{original equation}{original equations}
\Crefname{originalequation}{Original equation}{Original equations}
\makeatletter
\newlabel{equation-map-presentations}{{0H2F}{source}{}{sourcepointer-2}{}}
\newlabel{equation-map-presentations@cref}{{[originalequation][0][]0H2F}{[1][0][]source}{}{}{}}
\makeatother
\title{The map on differentials for a surjection, with its kernel identified}
\author{The Stacks Project Authors}
\date{stacks-project:algebra.tex:local-snapshot-20261012; GNU FDL 1.2 or later}
\begin{document}
\maketitle
\noindent\textit{Editorial note.} All delivered mathematical text is the author's own, copied from the pinned Stacks Project chapter file algebra.tex; no statement or proof step was rewritten, reordered or removed.\par\medskip
\noindent\textit{Editorial note.} Preamble provenance, as required for this source: the frozen snapshot stages the chapter's own preamble file beside it (preamble.tex, 217 lines, 8116 bytes, the same shared Stacks-wide preamble the published chow and etale-cohomology cards use); the first line of the pinned chapter file is \textbackslash{}input\{preamble\}, so no preamble text lives inside the chapter itself. The card takes that file's lines 8-18 (the verbatim package and the three \textbackslash{}comment-based environments reference/slogan/history at lines 9-12, section 13.2, together with the Xy-pic setup) and its lines 154 to the runtime end of the file (the native theorem declarations and mathematical macros). Not one character was changed, no package was added, no macro was defined, and the excerpt serves only as preamble, never as a span coordinate.\par\medskip
\noindent\textit{Editorial note.} The same-chapter pointers are located at their \textbackslash{}label lines inside the pinned chapter file; the cross-chapter pointers are located in the project's own tag list staged in the snapshot (sources/stacks-chow/tags.txt), at the verbatim lines carrying their canonical tags.\par\medskip
\noindent\textit{Editorial note.} The counted statement and proof refer to two numbered equations of the same chapter; they are recorded as navigation stubs with kind equation.\par\medskip
\noindent\textit{Editorial note.} The counted spans contain no \textbackslash{}cite, so the card carries no bibliography.\par\medskip
\noindent\textit{Allowed standard prerequisites (editorial).}\begin{itemize}
\item The map of presentations induced by a square of ring maps. Used at the start of the counted proof (lines 34204-34207). Gives the surjectivity of the map on differentials from the surjectivity of the top map of free modules.
\item Description of differentials of the target by generators and relations. Used in the kernel computation (lines 34208-34233). Lets an element mapping to zero be written in a normal form that exposes its coefficients.
\end{itemize}\medskip
\ifdefined\hypertarget\hypertarget{sourcepointer-1}{}\fi
\noindent\textit{Source reference.} Original equation 00RQ: sources/stacks-algebra/pinned-original-source.tex, source lines 34114--34114. Navigation pointer for \textbackslash{}ref\{equation-functorial-omega\} used inside the counted span (line 34196); the argument that consumes it is retained in the span.\par\medskip
\ifdefined\hypertarget\hypertarget{sourcepointer-2}{}\fi
\noindent\textit{Source reference.} Original equation 0H2F: sources/stacks-algebra/pinned-original-source.tex, source lines 34144--34144. Navigation pointer for \textbackslash{}ref\{equation-map-presentations\} used inside the counted span (line 34205); the argument that consumes it is retained in the span.\par\medskip
\begin{lemma}
\label{lemma-differential-surjective}
In diagram (\ref{equation-functorial-omega}), suppose
that $S \to S'$ is surjective with kernel $I \subset S$.
Then $\Omega_{S/R} \to \Omega_{S'/R'}$ is surjective with
kernel generated as an $S$-module by the elements
$\text{d}a$, where $a \in S$ is such that $\varphi(a) \in \beta(R')$.
(This includes in particular the elements $\text{d}(i)$, $i \in I$.)
\end{lemma}

\begin{proof}[First proof]
Consider the map of presentations (\ref{equation-map-presentations}).
Clearly the right vertical
map of free modules is surjective. Thus the map is surjective.
Suppose that some element $\eta$ of $\Omega_{S/R}$ maps to zero in
$\Omega_{S'/R'}$. Write $\eta$ as the image of $\sum s_i[a_i]$ for some
$s_i, a_i \in S$. Then we see that $\sum \varphi(s_i)[\varphi(a_i)]$
is the image of an element
$$
\theta = \sum s_j'[a_j', b_j'] + \sum s_k'[f_k', g_k'] + \sum s_l'[r_l']
$$
in the upper left corner of the diagram.
Since $\varphi$ is surjective, the terms
$s_j'[a_j', b_j']$ and $s_k'[f_k', g_k']$ are in the image
of elements in the lower right corner. Thus, modifying
$\eta$ and $\theta$ by subtracting the images of these
elements, we may assume $\theta = \sum s_l'[r_l']$.
In other words, we see $\sum \varphi(s_i)[\varphi(a_i)]$ is of the form
$\sum s'_l [\beta(r'_l)]$.
Next, we may assume that we have some $a' \in S'$
such that $a' = \varphi(a_i)$ for all $i$ and $a' = \beta(r_l')$
for all $l$. This is clear from the direct sum decomposition
of the upper right corner of the diagram. Choose $a \in S$ with
$\varphi(a) = a'$. Then we can write $a_i = a + x_i$
for some $x_i \in I$. Thus we may assume that all $a_i$ are equal to
$a$ by using the relations that are allowed. But then we may
assume our element is of the form $s[a]$. We still know that
$\varphi(s)[a'] = \sum \varphi(s_l')[\beta(r_l')]$.
Hence either $\varphi(s) = 0$ and we're done, or
$a' = \varphi(a)$ is in the image of $\beta$ and we're done as well.
\end{proof}

\end{document}
